An equational theory for a family of quantum circuits is complete if it
proves every equation that holds between the matrices of the circuits.
Completeness proofs are long and repetitive. They normalize circuits, or
synthesize them from their matrices, and then check by hand that every
rule meets every normal form in the expected way. We have checked twelve
such proofs in Agda, written between 2014 and 2026, together with the
minimality theorem of a thirteenth paper. They are about Clifford, real
stabilizer, CNOT-dihedral, phase-affine, Clifford+$T$, Clifford+$CS$ and
real Clifford+$CH$ circuits, and about the groups of unitary and
orthogonal matrices over $\Z[\tfrac12,i]$, $\Z[\tfrac1{\sqrt2},i]$,
$\Z[\tfrac12]$ and $\Z[\tfrac1{\sqrt2}]$ in which such circuits live.
The development has about 950 modules and 250{,}000 lines of safe Agda
without postulates. It was written in 24 days with an LLM coding agent.

Most of the proofs are correct, up to misprints that we list. One is
not. The completeness proof for Figure~6 of Hadamard-$\Pi$ (POPL 2026)
relies on two false lemmas, and its Main Lemma fails on an explicit state
in dimension~9. The theorem itself is true. We prove it by formalizing
the corresponding case of an unpublished 2019 thesis of Cl\'ement, and
then the whole thesis. Cl\'ement's 2026 completeness theorem for real
Clifford+$CH$ circuits, whose appendix has 264 pages,
imports this theorem; our proof of it has no such hypothesis. We
describe the proof architectures the papers share, how we found and
checked hundreds of derivations, and what the work cost in Agda.
